We formalize QRC and ESN forecasting as fixed feature-map comparisons with a trained readout. For a univariate or input-driven time series, a reservoir maps the history up to time \(t\) into a feature vector \(\Phi_t\in\R^D\). The readout is fit on a training split and evaluated on a held-out test split. The primary error metric is normalized root mean squared error,
\begin{equation}
    \nrmse =
    \frac{\sqrt{\frac{1}{T}\sum_t (y_t-\hat y_t)^2}}
    {\operatorname{std}(y_t)} .
    \label{eq:nrmse}
\end{equation}

\subsection{QRC as a fixed measured feature map}

For QRC, the feature map is generated by fixed quantum dynamics. Given input \(u_t\), the reservoir state evolves under an input-dependent quantum channel
\begin{equation}
    \rho_t = \mathcal{Q}_{u_t}(\rho_{t-1}),
    \qquad
    \Phi_t =
    \big(\Tr[O_a\rho_t]\big)_{a=1}^{D},
    \label{eq:qrc-feature}
\end{equation}
where \(O_a\) are measured observables. Only the classical readout is trained. In the v5 protocol, the QRC family consists of three canonical frozen variants: a mechanistic disordered spin-QRC (QRC-M), an encoding-enhanced variant with fixed reuploading (QRC-E), and a mildly dissipative variant (QRC-D). All use the same feature dimension as the ESN comparator.

\subsection{ESN comparator}

Echo state networks are the classical reservoir-computing baseline most directly matched to QRC: random recurrent dynamics produce a nonlinear state, and only the readout is trained~\cite{jaeger2004harnessing,lukosevicius2009reservoir}. Our ESN is a sparse random leaky reservoir with globally calibrated spectral radius, leak rate, and input scale. Its reservoir size is matched to the QRC feature dimension. This makes the comparison a fixed quantum feature map versus a fixed classical recurrent feature map, not a per-dataset tuning contest.

\subsection{Dataset-regime view}

The atlas treats QRC advantage as a property of a dataset regime. For each dataset \(i\), define
\begin{equation}
    A_i =
    \nrmse_{\esn,i}
    -
    \min_{v\in\{\mathrm{M,E,D}\}}\nrmse_{\qrc_v,i}.
    \label{eq:advantage}
\end{equation}
Positive \(A_i\) means that the best frozen QRC variant beats the frozen ESN. We call a case \emph{useful} when \(A_i\geq 0.05\). The meta-model then asks whether measured descriptors \(z_i\) predict the label \(A_i\geq 0.05\) before running the QRC.
